\chapter{Conclusions and Future Work}

This work presented a comprehensive study on the application of Multi-Agent Reinforcement Learning (MARL)
approaches to the problem of underwater search and navigation in an unknown, realistic and complex environment.
The study focused on the application of multi-agent PPO variants, especially IPPO and MAPPO, with the addition of a 
communication mechanism to enhance the agents' performance in a partially observable environment. The results demonstrated
the effectiveness of the proposed approaches in making the agents find the target area in the shortest amount of time. 
Different approaches lead to different results, even though independent agents are optimal choice if there is not prior 
information regarding the environment. In fact, the use of communication has been proved to help in difficult scenarios and
in situations where agents are spread and thus find helpful having knowledge about other points in the environment,
whether if they start close it does not help much. The centralized critic was not generally found helpful in this task,
even though in some cases it helped achieving a better performance, especially with a lower number of agents and with a 
smaller target area. 
In this problem, since the target is defined as an area and usually it is quite large, training a policy with more than 2 agents
was difficult and not optimal. Increasing to 4 agents, in fact, showed that independence leads to better performances overall, which is 
different from what one might expect. The effect of lucky spawning near the target zone, more probable in this case,
has led to an optimal policy which favors the removal of communication, probably because using a larger state does not bring 
any advance with respect to the independent case, which is comprehensible.
Efficiency is measured as the time to reach success, and in this task independent agents reached the worst performance overall, 
while MAPPO is the most efficient policy with 2 agents. This is explained by the fact that agents trained with a centralized policy
are able to divide the work more efficiently than independent agents, which only care about success and reaching target. Coordinated agents,
instead, trade some success for more efficiency, and that is proved by the time they take to reach the target zone. 
Finally, the way agents are initialized has been found to matter a lot,
and the verdict is to use a uniform initialization to encode a natural spread in the environment. This could be informative for 
real deployment consideration, since it is not said that the cost of releasing the AUVs spread is lower than the what the difference in 
success rate brings, and that should be defined beforehand.

This problem is inherently difficult, since the target area depends on how the environment is made up and which variables
are considered. In fact, in this case only two parameters are considered for the target, but there could be more. Adding more 
constraints naturally complicates the learning problem, and there could be no optimal policy to navigate the environment in that case,
since the more the variables that need to be taken into consideration, the more difficult would be for the agent to navigate using gradients
and the higher will be the dimension of the state space. However, this should be the very next starting point for further analyze the effect 
of scaling the number of agents, since the dimension of the target area resulted to be the main problem that limited the 
training of an optimal policy with more than 2 agents. Moreover, 
other forms of communication and cooperation could be exploited and tested, for example with swarm capabilities that in this case were unhelpful due to the 
small environment and number of agents considered. If a larger number of agents are used, which complicates even further the learning 
problem, the consideration of techniques like PSO could help the agents avoid local optima generated from the more complex environment,
since those techniques are more effective with large swarms. 

To conclude, learning optimal policies in this task is definitely possible, even though the complexity of the real environment 
may be much higher than a simulated one like the one developed in this work. In fact, obstacles may be present and unexpected hazards
are possible and should be taken into account. This work proved that MARL is a promising approach to solve problems in the underwater domain,
like coral reef restoration and monitoring, which are vital for the health of the marine ecosystem. Further optimizations and improvements are needed 
to make the proposed approaches more robust and effective in real-world scenarios, but the results obtained in this work provide a solid foundation for future research in this area.